% ═══════════════════════════════════════════════════════════════════
%  INFORME — Tarea IC RNA 2026-1
%  Paper: Widrow & Lehr (1990)
%  "30 Years of Adaptive Neural Networks:
%   Perceptron, Madaline, and Backpropagation"
% ═══════════════════════════════════════════════════════════════════

\documentclass[11pt,a4paper]{article}

% ── Paquetes ──────────────────────────────────────────────────────
\usepackage[utf8]{inputenc}
\usepackage[T1]{fontenc}
\usepackage[spanish,es-noshorthands]{babel}
\usepackage{amsmath,amssymb,amsfonts}
\usepackage{graphicx}
\usepackage[margin=2.3cm]{geometry}
\usepackage{float}
\usepackage{booktabs}
\usepackage{enumitem}
\usepackage{caption}
\usepackage[colorlinks=true,linkcolor=blue,citecolor=blue,urlcolor=blue]{hyperref}

% ── Macros para figuras y tablas (compilacion robusta) ───────────
\newcommand{\figura}[4][0.92]{%
\begin{figure}[H]\centering
\IfFileExists{../../notebook/v1/figs/#2.png}{%
\includegraphics[width=#1\textwidth]{../../notebook/v1/figs/#2.png}}{%
\fbox{\parbox[c][2.6cm][c]{0.8\textwidth}{\centering\texttt{figs/#2.png}\\(pendiente: ejecutar el notebook)}}}
\caption{#3}\label{fig:#4}
\end{figure}}

\newcommand{\tabla}[3]{%
\begin{table}[H]\centering\small
\IfFileExists{../../notebook/v1/tabs/#1.tex}{%
\input{../../notebook/v1/tabs/#1.tex}}{%
\fbox{\parbox[c][1.6cm][c]{0.7\textwidth}{\centering\texttt{tabs/#1.tex}\\(pendiente)}}}
\caption{#2}\label{tab:#3}
\end{table}}

% ── Comandos personalizados ──────────────────────────────────────
\newcommand{\R}{\mathbb{R}}
\newcommand{\E}{\mathbb{E}}
\newcommand{\norm}[1]{\lVert#1\rVert}
\newcommand{\aLMS}{$\alpha$-LMS}
\newcommand{\muLMS}{$\mu$-LMS}
\newcommand{\bp}{\textit{backpropagation}}
\newcommand{\mrii}{\textit{Madaline Rule II}}
\newcommand{\mriii}{\textit{Madaline Rule III}}

% ── Metadatos ─────────────────────────────────────────────────────
\title{\textbf{Algoritmos de Aprendizaje en Redes Neuronales Feedforward}\\[4pt]
\large Un análisis del artículo de Widrow \& Lehr (1990):\\
\textit{30 Years of Adaptive Neural Networks: Perceptron, Madaline, and Backpropagation}}
\author{Estudiante\\Universidad de Santiago de Chile\\Magíster en Ingeniería Informática}
\date{Julio 2026}

\begin{document}
\maketitle

\begin{abstract}
Este informe presenta un análisis comprehensivo del artículo de revisión de Widrow \& Lehr (1990)
sobre 30 años de algoritmos de aprendizaje supervisado en redes neuronales artificiales (RNA)
feedforward. Se implementaron desde cero en NumPy puro los seis algoritmos descritos en el artículo
—Perceptrón, \aLMS, \muLMS, MLP con \bp, \mrii~y \mriii— y se ejecutaron siete experimentos
para verificar experimentalmente las hipótesis fundamentales del principio de mínima perturbación.
Los resultados confirman que: (1) la capacidad estadística $C_s \approx 2N_w$ de Cover se verifica
empíricamente, (2) el Perceptrón converge en tiempo finito solo en datos separables mientras \muLMS
se aproxima al óptimo de Wiener, (3) el MLP con \bp~resuelve XOR (6/20, 2--2--1) y alcanza 91.1\% en
Iris, (4) \mriii~es matemáticamente equivalente a \bp~con error relativo de $5.05 \times 10^{-8}$
pero entre 3.5$\times$ y 74$\times$ más costoso computacionalmente. Se discuten las implicancias
del principio de mínima perturbación como marco unificador del aprendizaje supervisado en RNA.
\end{abstract}

% ═══════════════════════════════════════════════════════════════════
%  1. INTRODUCCIÓN (15%)
% ═══════════════════════════════════════════════════════════════════
\section{Introducción}

\subsection{Problema abordado}

El artículo de Widrow \& Lehr (1990) \cite{widrow1990} aborda el problema fundamental del
\textbf{aprendizaje supervisado en redes neuronales artificiales feedforward}: ¿cómo adaptar los
pesos sinápticos de una red de elementos computacionales simples para que aproximen una función
de clasificación o regresión a partir de ejemplos? El artículo revisa 30 años de desarrollo de
algoritmos de entrenamiento —desde el Perceptrón de Rosenblatt (1958) hasta la retropropagación
del error de Rumelhart, Hinton \& Williams (1986)— y demuestra que todos ellos comparten un
principio unificador: el \textbf{principio de mínima perturbación} (\textit{minimal disturbance
principle}), que establece que durante el entrenamiento es recomendable inyectar nueva información
en la red de manera que perturbe la información ya almacenada en la menor medida posible
\cite{widrow1990}.

El artículo es una revisión (\textit{survey}) y tutorial, no un estudio experimental sobre un
dataset único. Su contribución es \textbf{conceptual e histórica}: traza la evolución desde las
primeras reglas de aprendizaje (1960) hasta las redes multicapa con \bp, y demuestra conexiones
matemáticas profundas entre algoritmos desarrollados independientemente. En particular, establece
que \mriii~(Andes, 1988) es matemáticamente equivalente a \bp, un hallazgo realizado por Widrow
y sus estudiantes \cite{widrow1990}.

\subsection{¿Por qué redes neuronales artificiales?}

La motivación para usar RNA en problemas de clasificación y regresión se fundamenta en tres
argumentos presentados en las Secciones I--II del artículo \cite{widrow1990}:

\begin{enumerate}[label=\textbf{(\roman*)},leftmargin=*]
    \item \textbf{Paralelismo masivo y robustez:} Los sistemas biológicos resuelven tareas de
    reconocimiento de patrones, procesamiento de voz e imágenes con una facilidad que las
    computadoras von Neumann no igualan. Las RNA emulan este paralelismo mediante miles de
    elementos simples interconectados que procesan información simultáneamente.
    \item \textbf{Aprendizaje por ejemplos:} A diferencia de los sistemas algorítmicos
    tradicionales, las RNA no requieren programación explícita de reglas. Aprenden la función
    de mapeo entrada-salida directamente de los datos mediante algoritmos iterativos.
    \item \textbf{Generalización:} Una RNA correctamente entrenada clasifica patrones no vistos
    durante el entrenamiento. La capacidad de generalización está inversamente relacionada con la
    capacidad de almacenamiento ($C_s \approx 2N_w$, Cover, 1964 \cite{cover1965}), un
    \textit{trade-off} fundamental que el artículo explora en profundidad.
\end{enumerate}

Aplicaciones históricas que demuestran la viabilidad práctica incluyen NETtalk (Sejnowski \&
Rosenberg, 1987), una red de 80--26 Adalines que aprende a pronunciar texto en inglés; el
\textit{truck backer-upper} de Nguyen \& Widrow (1989), donde una red de 26 Adalines aprende a
retroceder un camión con acoplado; y los ecualizadores adaptativos basados en LMS para módems
de alta velocidad, la primera aplicación comercial masiva de las RNA \cite{widrow1990}.

\subsection{Hipótesis}

Se plantean cuatro hipótesis que guían la experimentación de este informe:

\begin{description}
    \item[H1 -- Capacidad de Cover:] La capacidad estadística de un Adaline es $C_s \approx 2N_w$
    y la determinística $C_d = N_w$, verificable mediante simulación Monte Carlo con programación
    lineal.
    \item[H2 -- Convergencia finita del Perceptrón:] El Perceptrón converge en tiempo finito si
    los datos son linealmente separables, pero oscila y diverge en norma si no lo son. Por
    contraste, \aLMS~y \muLMS~son estables en ambos casos, con \muLMS~convergiendo al óptimo de
    Wiener $W^* = R^{-1}P$.
    \item[H3 -- Poder no lineal y óptimos locales:] XOR es irresoluble por un solo Adaline pero
    resoluble por \mrii~y MLP con \bp, aunque ambos algoritmos están sujetos a óptimos locales
    en la superficie de error. El remedio de ruido esporádico propuesto en la Sección VII-E del
    paper \cite{widrow1990} mitiga parcialmente este problema.
    \item[H4 -- Equivalencia MRIII $\equiv$ Backpropagation:] \mriii~converge a \bp~cuando la
    perturbación $\Delta s \to 0$, pero su costo computacional en implementación digital es
    proporcionalmente mayor (requiere $1 + N_{\text{Adalines}}$ pasadas hacia adelante por
    presentación versus 1 adelante $+$ 1 atrás de \bp).
\end{description}

% ═══════════════════════════════════════════════════════════════════
%  2. MARCO TEÓRICO (15%)
% ═══════════════════════════════════════════════════════════════════
\section{Marco Teórico}

\subsection{El elemento adaptativo básico: Adaline}

El bloque constructivo fundamental de las RNA feedforward es el \textbf{Adaptive Linear Element}
(Adaline) \cite{widrow1990}, compuesto por:

\begin{enumerate}[leftmargin=*]
    \item Un \textbf{combinador lineal adaptativo} que calcula $s_k = X_k^T W_k$, donde
    $X_k = [x_0=+1, x_{1k}, \ldots, x_{nk}]^T$ es el vector de entrada aumentado con bias
    y $W_k \in \R^{n+1}$ el vector de pesos.
    \item Un \textbf{cuantizador} (típicamente signum) que produce la salida binaria
    $y_k = \mathrm{sgn}(s_k) \in \{\pm1\}$, o alternativamente una función de activación
    diferenciable como la tangente hiperbólica $y_k = \tanh(s_k)$.
\end{enumerate}

La condición crítica $s = 0$ define un \textbf{hiperplano separador} en el espacio de patrones.
Con $n$ entradas binarias, un Adaline puede implementar solo las funciones lógicas
\textbf{linealmente separables}: 14 de 16 posibles para $n=2$, excluyendo XOR/XNOR \cite{widrow1990}.

\subsection{Taxonomía de algoritmos}

El artículo organiza los algoritmos en una taxonomía (Fig.~33 del paper) según dos ejes:
\textbf{mecanismo de adaptación} (corrección de error vs. descenso más pronunciado) y
\textbf{no linealidad} (lineal vs. no lineal), para uno o múltiples elementos:

\begin{enumerate}[leftmargin=*]
    \item \textbf{Corrección de error -- lineal (un elemento):} \aLMS~(Widrow-Hoff, 1960),
    que minimiza el error lineal $\varepsilon_k = d_k - s_k$ con cambio de peso colineal a $X_k$.
    \item \textbf{Corrección de error -- no lineal (un elemento):} Regla del Perceptrón
    (Rosenblatt, 1958), que solo adapta si la decisión es incorrecta usando el error del
    cuantizador $\tilde{\varepsilon}_k = d_k - y_k$. Reglas de Mays con zona muerta.
    \item \textbf{Corrección de error -- no lineal (red):} Madaline I (MRI) y II (MRII):
    MRI usa capa oculta adaptativa con lógica fija de salida; MRII extiende la adaptación
    a todas las capas mediante inversión tentativa (\textit{tentative flipping}) \cite{widrow1990}.
    \item \textbf{Descenso más pronunciado -- lineal (un elemento):} \muLMS, que realiza
    descenso por gradiente estocástico sobre la superficie de MSE, la cual es un paraboloide
    convexo con mínimo único: la solución de Wiener $W^* = R^{-1}P$.
    \item \textbf{Descenso más pronunciado -- no lineal (red):} \bp~(Rumelhart et al., 1986)
    y \mriii~(Andes, 1988), que extienden el gradiente a redes multicapa con funciones de
    activación diferenciables.
\end{enumerate}

\subsection{Ecuaciones fundamentales}

\subsubsection{$\alpha$-LMS / Regla Delta de Widrow-Hoff — Eq.~(10)}

\begin{equation}
W_{k+1} = W_k + \alpha\,\frac{\varepsilon_k}{\norm{X_k}^2}\,X_k,\quad \varepsilon_k = d_k - X_k^T W_k
\label{eq:alpha_lms}
\end{equation}

El cambio de peso $\Delta W_k$ es colineal con $X_k$, garantizando la corrección del error con
la mínima magnitud de cambio posible (\textit{mínima perturbación geométrica}). Rango de
estabilidad: $0 < \alpha < 2$, rango práctico: $0.1 < \alpha < 1.0$ \cite{widrow1990}.

\subsubsection{Regla del Perceptrón de Rosenblatt — Eq.~(18)}

\begin{equation}
W_{k+1} = W_k + \alpha\,\frac{\tilde{\varepsilon}_k}{2}\,X_k,\quad \tilde{\varepsilon}_k = d_k - y_k \in \{-2,0,+2\}
\label{eq:perceptron}
\end{equation}

Solo adapta si $y_k \neq d_k$. El \textbf{teorema de convergencia del Perceptrón} garantiza
separación en pasos finitos si los datos son linealmente separables. Si no lo son, el vector
de pesos tiende a cero y el algoritmo ``continúa para siempre'' \cite{widrow1990,rosenblatt1958}.

\subsubsection{$\mu$-LMS — Eq.~(33)}

\begin{equation}
W_{k+1} = W_k + 2\mu\,\varepsilon_k\,X_k
\label{eq:mu_lms}
\end{equation}

La superficie de MSE $\xi(W) = \E[d^2] - 2P^T W + W^T R W$ es un hiperparaboloide convexo
con mínimo único $W^* = R^{-1}P$ (solución de Wiener). El gradiente instantáneo es insesgado:
$\E[-2\varepsilon_k X_k] = \nabla\xi$. Estabilidad: $0 < \mu < 1/\mathrm{tr}[R]$ \cite{widrow1990}.

\subsubsection{Backpropagation — Eqs.~(78--79, 87, 94)}

Para el Adaline sigmoide ($y = \tanh(s)$), la regla de actualización es:

\begin{equation}
W_{k+1} = W_k + 2\mu\,\tilde{\varepsilon}_k \cdot (1 - y_k^2)\,X_k
\label{eq:backprop_single}
\end{equation}

Para redes multicapa, las deltas se calculan como:
\begin{align}
\delta^{(L)}_j &= \varepsilon_j \cdot (1 - y_j^2) \quad \text{(capa de salida, Eq.~79)} \label{eq:delta_out}\\
\delta^{(l)} &= \bigl(W^{(l+1)T} \delta^{(l+1)}\bigr) \odot (1 - (y^{(l)})^2) \quad \text{(capas ocultas, Eq.~87)} \label{eq:delta_hid}
\end{align}

Y la actualización con \textbf{momentum} (Eqs.~96--97):
\begin{equation}
\Delta W_k = \eta\,\Delta W_{k-1} + (1-\eta)\,2\mu\,\delta_k X_k^T,\quad \eta \approx 0.8\text{--}0.9
\label{eq:momentum}
\end{equation}

Es crucial aclarar que $\delta$ \textbf{no es un error} en el sentido convencional: es una
señal de gradiente que indica cómo cambiar cada peso para reducir el error cuadrático de salida
$\varepsilon^2$ \cite{widrow1990,rumelhart1986}.

\subsection{Capacidad de patrones y generalización}

Cover (1964) \cite{cover1965} demostró que la probabilidad de que $N_p$ patrones aleatorios sean
linealmente separables por un Adaline de $N_w$ pesos es:

\begin{equation}
P(N_p, N_w) =
\begin{cases}
1 & \text{si } N_p \leq N_w \\[4pt]
2^{1-N_p} \displaystyle\sum_{i=0}^{N_w-1} \binom{N_p-1}{i} & \text{si } N_p > N_w
\end{cases}
\label{eq:cover}
\end{equation}

De aquí se derivan la \textbf{capacidad determinística} $C_d = N_w$ (garantía de solución) y la
\textbf{capacidad estadística} $C_s \approx 2N_w$ (promedio de patrones absorbibles). Para redes
multicapa, Baum extendió estos resultados acotando $N_w/N_y - K_1 \leq C_d \leq
(N_w/N_y)\log_2(N_w/N_y) + K_2$ \cite{widrow1990}. La regla práctica de generalización es
$N_p \gg N_w/N_y$: el número de patrones de entrenamiento debe exceder significativamente la
relación pesos/salidas para que la red generalice en lugar de memorizar.

\subsection{Madaline Rule II (MRII)}

\mrii~(Widrow, Winter \& Baxter, 1987) \cite{widrow1990} extiende la adaptación a todas las capas
de una red de Adalines con cuantizador signum. Para cada patrón erróneo:

\begin{enumerate}[leftmargin=*]
    \item Se ordenan las Adalines ocultas por $|s|$ creciente (\textit{load-sharing}: ``asignar
    responsabilidad a quien puede asumirla más fácilmente'').
    \item Se invierte tentativamente (\textit{flip}) la salida de una Adaline; si reduce el error
    de Hamming del patrón, se acepta y se refuerza con corrección absoluta colineal.
    \item Si no, se prueban pares de Adalines.
    \item La capa de salida se entrena con \aLMS~moderado.
\end{enumerate}

El artículo advierte que MRII puede estancarse en óptimos locales, y propone como remedio la
\textbf{adición esporádica de ruido a los pesos} (Sección VII-E) \cite{widrow1990}.

\subsection{Madaline Rule III (MRIII) y equivalencia con Backpropagation}

\mriii~(Andes, 1988) reemplaza el cuantizador signum por sigmoides y estima el gradiente de
$\varepsilon^2$ respecto a $s_j$ mediante perturbación del sumidero \cite{widrow1990}:

\begin{equation}
\frac{\partial \varepsilon^2}{\partial s_j} \approx
\frac{\varepsilon^2_{\text{perturbado}} - \varepsilon^2_{\text{base}}}{\Delta s}
\label{eq:mriii}
\end{equation}

Cuando $\Delta s \to 0$, MRIII es \textbf{matemáticamente equivalente a backpropagation}.
Requiere $1 + N_{\text{Adalines}}$ pasadas hacia adelante por patrón (vs. $1+1$ de \bp),
lo que lo hace menos eficiente en implementación digital, pero más robusto para hardware
analógico (existió un chip comercial de MRIII) \cite{widrow1990}.

% ═══════════════════════════════════════════════════════════════════
%  3. EXPERIMENTACIÓN (15%)
% ═══════════════════════════════════════════════════════════════════
\section{Experimentación}

\subsection{Diseño experimental}

Dado que el artículo de Widrow \& Lehr es una revisión teórica, la experimentación propia
consiste en \textbf{implementar desde cero los algoritmos descritos} y verificar empíricamente
las hipótesis H1--H4 mediante siete experimentos (E1--E7), más validación en dos datasets reales
adicionales (Wine y MNIST reducido). La Tabla~\ref{tab:diseno} resume el diseño experimental.

\begin{table}[H]\centering\small
\begin{tabular}{@{}cllll@{}}
\toprule
\textbf{Exp.} & \textbf{Tema} & \textbf{Hipótesis} & \textbf{Algoritmos} & \textbf{Dataset} \\
\midrule
E1 & Capacidad de Cover & H1 & Programación Lineal & Sintético Monte Carlo \\
E2 & Reglas lineales, separable & H2 & Perceptrón, \aLMS, \muLMS & Gaussianas separables \\
E3 & Datos no separables & H2 & Perceptrón, \aLMS, \muLMS & Gaussianas solapadas \\
E4 & Superficies MSE & — & Adaline 2D & Patrones aleatorios \\
E5 & XOR: MRII vs MLP & H3 & MRII, MLP-\bp & XOR (4 patrones) \\
E6 & Datasets reales & — & MLP-\bp, \aLMS, \muLMS & Iris, Wine, MNIST-04, two-moons \\
E7 & MRIII $\equiv$ \bp~+ costo & H4 & MRIII, \bp & Iris (1 patrón) \\
\bottomrule
\end{tabular}
\caption{Diseño experimental.}\label{tab:diseno}
\end{table}

\subsection{Datasets}

\begin{itemize}[leftmargin=*]
    \item \textbf{Sintéticos separables:} 200 puntos, dos gaussianas $N([\pm2,\pm2], 0.35^2I)$.
    \item \textbf{Sintéticos no separables:} 200 puntos, dos gaussianas $N([\pm0.9,\pm0.9], 1.3^2I)$.
    \item \textbf{XOR:} 4 patrones $\{\pm1\}^2$, $d = [-1, +1, +1, -1]$.
    \item \textbf{Iris} (Fisher, 1936; UCI): 150 flores, 4 atributos, 3 clases. Split 70/30 estratificado, normalizado a $[-1,1]$.
    \item \textbf{Two-moons} (scikit-learn): 300 puntos, ruido 0.20.
    \item \textbf{Wine} (UCI): 178 muestras, 13 atributos fisicoquímicos, binarizado (clase 0 vs. resto).
    \item \textbf{MNIST-04} (UCI digits): dígitos 0--4, 8$\times$8=64 atributos, 901 instancias, normalizado.
\end{itemize}

\subsection{Implementación}

Todos los algoritmos se implementaron desde cero en \textbf{NumPy puro} (sin usar PyTorch,
TensorFlow ni scikit-learn para el aprendizaje). Se usó scikit-learn exclusivamente para carga
de datos, partición train/test y métricas auxiliares. La reproducibilidad se garantiza con
semilla global \texttt{SEED=1990}. El costo computacional por presentación es $\mathcal{O}(N_w)$
para Perceptrón/LMS y $\mathcal{O}(N_w)$ para \bp~(forward $+$ backward). Para MRIII, el costo
es $\mathcal{O}(N_w \cdot N_{\text{Adalines}})$.

\subsection{Complejidad computacional}

En implementación digital, \bp~requiere exactamente 1 pasada hacia adelante $+$ 1 hacia atrás por
patrón, con complejidad comparable en ambas direcciones. MRIII requiere $1 + N_{\text{Adalines}}$
pasadas hacia adelante (una sin perturbar, una por cada Adaline perturbada), lo que lo hace
intrínsecamente más costoso. Esta diferencia se cuantifica en el Experimento E7.

% ═══════════════════════════════════════════════════════════════════
%  4. ANÁLISIS DE RESULTADOS (30%)
% ═══════════════════════════════════════════════════════════════════
\section{Análisis de Resultados}

\subsection{E1: Capacidad de Cover}

\figura{fig_cover}{Probabilidad de separabilidad lineal: empírica (puntos) vs.~fórmula teórica de Cover (curvas). Líneas verticales en $C_d=N_w$ (punteada) y $C_s \approx 2N_w$ (discontinua).}{cover}

La Figura~\ref{fig:cover} muestra un ajuste estrecho entre la probabilidad empírica (80 ensayos
Monte Carlo por punto, separabilidad testeada con \texttt{linprog}) y la fórmula teórica de
Cover (Eq.~\ref{eq:cover}). Las desviaciones medias $|\text{empírica} - \text{teórica}|$ son
$0.0337$ para $N_w=2$, $0.0259$ para $N_w=5$ y $0.0105$ para $N_w=15$, decreciendo con $N_w$
como predice la teoría asintótica. La transición en $N_p/N_w = 1$ ($C_d$) y $N_p/N_w = 2$
($C_s$) se agudiza al aumentar $N_w$, confirmando \textbf{H1}.

\subsection{E2: Reglas lineales en datos separables}

\figura{fig_datasets}{Visualización de los datasets sintéticos utilizados: separable,
no separable, y XOR.}{datasets}

\figura{fig_fronteras}{Fronteras de decisión de los tres algoritmos lineales (color) vs.~el
hiperplano óptimo de Wiener (negro discontinuo). Ángulos respecto a $W^*$ indicados.}{fronteras}

\figura{fig_convergencia}{Convergencia de \muLMS~para distintos $\mu$. $\mu=0.12$ diverge por
exceder la cota de estabilidad $\mu_{\max} = 1/\mathrm{tr}[R] = 0.111$. La línea horizontal
marca el MSE del óptimo de Wiener (0.015).}{convergencia}

La Figura~\ref{fig:fronteras} revela un resultado fundamental: \textbf{el Perceptrón converge
en solo 2 épocas} (confirmando el teorema de convergencia), pero su frontera queda a $20.6^\circ$
del óptimo de Wiener — el Perceptrón ``se detiene donde primero separa''. En contraste, \aLMS~y
\muLMS~optimizan el MSE, acercándose a $4.4^\circ$ y $3.5^\circ$ de $W^*$ respectivamente.

La Figura~\ref{fig:convergencia} demuestra la cota de estabilidad: con $\mu=0.12 >
1/\mathrm{tr}[R] = 1/9.02 = 0.111$, el MSE diverge exponencialmente. Los valores $\mu \in
\{0.001, 0.01, 0.05\}$ convergen al MSE de Wiener ($0.015$), con mayor $\mu$ implicando
convergencia más rápida. La \textbf{MSE de Wiener} ($0.0148$) se alcanza asintóticamente.
Estos resultados confirman \textbf{H2} para el caso separable.

\subsection{E3: Datos no separables}

\figura{fig_noseparable}{Izquierda: Accuracy por época en datos no separables. Derecha: Evolución
de $\norm{W}$. Nótese la oscilación del Perceptrón (azul).}{noseparable}

La Tabla~\ref{tab:lineales} resume los resultados cuantitativos.

\tabla{tab_lineales}{Resultados comparativos de algoritmos lineales en datos separables y no separables.}{lineales}

En el dataset no separable (Figura~\ref{fig:noseparable}), el \textbf{Perceptrón oscila}
significativamente: $75.7\% \pm 5.5$ en las últimas 40 épocas, con $\norm{W}$ fluctuante (el
vector de pesos ``gravita hacia cero'' como advierte el paper). \aLMS~es estable pero ruidoso
($75.0\% \pm 7.4$). \muLMS~ofrece el mejor desempeño: $79.6\% \pm 1.5$, a solo 1.4 puntos
porcentuales del óptimo de Wiener ($81.0\%$), y con la menor varianza. Esto confirma
\textbf{H2} en el caso no separable: \muLMS~proporciona una solución estable cercana al
mejor clasificador lineal posible.

\subsection{E4: Superficies de MSE}

\figura{fig_superficies_mse}{Superficies de MSE para un Adaline de 2 pesos sin bias
(6 patrones aleatorios). Izquierda: error lineal (paraboloide convexo). Centro: error
sigmoide (no cuadrático pero suave). Derecha: error signum (escalonado con mesetas
y óptimos locales).}{superficies}

La Figura~\ref{fig:superficies} ilustra por qué el descenso de gradiente \textbf{exige funciones
de activación diferenciables}:

\begin{itemize}[leftmargin=*]
    \item El error \textbf{lineal} ($\min = 0.936$) forma un paraboloide perfectamente convexo:
    cualquier algoritmo de gradiente converge al mínimo global.
    \item El error \textbf{sigmoide} ($\min = 0.940$) es no cuadrático pero suave y navegable:
    el gradiente existe en todas partes, aunque pueden existir mínimos locales en redes multicapa.
    \item El error \textbf{signum} ($\min = 1.333$) es escalonado: presenta \textbf{mesetas}
    donde el gradiente es exactamente cero y la optimización se detiene. Esto explica por qué
    algoritmos como MRII (que usan signum) requieren mecanismos de escape como el ruido
    esporádico.
\end{itemize}

\subsection{E5: XOR — MRII vs MLP-Backpropagation}

\figura{fig_xor}{Izquierda: Regiones de decisión de MRII 2--3--1 con ruido. Centro: MLP-\bp~2--2--1.
Derecha: Curva de aprendizaje MSE del MLP.}{xor}

\tabla{tab_xor}{Tasa de éxito en XOR sobre 20 inicializaciones distintas.}{xor}

La Tabla~\ref{tab:xor} presenta el resultado más revelador de este experimento:

\begin{itemize}[leftmargin=*]
    \item \textbf{MRII 2--2--1 sin ruido: 0/20.} Sin el mecanismo de escape, MRII se estanca
    sistemáticamente en óptimos locales, tal como advierte la Sección VII-E del paper. Este es
    un \textbf{resultado}, no un error de implementación.
    \item \textbf{MRII con ruido esporádico ($\sigma=0.15$): 0/20 incluso con 3 Adalines ocultas.}
    En nuestra implementación, el ruido gaussiano no fue suficiente para escapar de los mínimos
    locales en XOR, posiblemente porque la corrección absoluta agresiva del refuerzo colineal
    domina sobre el escape estocástico.
    \item \textbf{MLP-\bp~2--2--1: 6/20} (mediana $\approx$ 40 épocas). La inicialización
    aleatoria de pesos determina si la red cae o no en un mínimo local, consistente con la
    literatura \cite{widrow1990,rumelhart1986}.
    \item \textbf{MLP-\bp~2--3--1: 11/20.} Una Adaline adicional en la capa oculta casi duplica
    la tasa de éxito al proporcionar más grados de libertad para sortear los mínimos locales.
\end{itemize}

Estos resultados confirman \textbf{H3}: XOR es resoluble por arquitecturas multicapa pero tanto
MRII como \bp~están sujetos a óptimos locales. La diferencia en las tasas de éxito (6/20 para
MLP vs. 0/20 para MRII) refleja que la superficie de error sigmoide (E4) es inherentemente más
navegable que la del signum.

\subsection{E6: MLP en datasets reales}

\figura{fig_iris_mlp}{Resultados del MLP 4--8--3 en Iris. Izquierda: Curva MSE. Centro: Accuracy
por época. Derecha: Matriz de confusión (91.1\% test).}{iris_mlp}

\figura{fig_moons}{Regiones de decisión del MLP 2--8--1 en two-moons (93.3\% test).}{moons}

\figura{fig_wine}{Convergencia de clasificadores lineales en Wine. Wiener y \muLMS~alcanzan
98.1\% test.}{wine}

\figura{fig_mnist}{MLP 64--16--5 en MNIST dígitos 0--4 (98.5\% test).}{mnist}

\tabla{tab_iris}{Resultados del MLP 4--8--3 en Iris para la grilla $\mu \times \eta$.}{iris}

\tabla{tab_relevancia}{Relevancia de variables en Iris medida como $\norm{W_1}_2$ por atributo.}{relevancia}

\textbf{Iris (Figura~\ref{fig:iris_mlp}, Tablas~\ref{tab:iris}--\ref{tab:relevancia}):} Como
predice la teoría, Setosa es linealmente separable del resto: el Perceptrón converge en 2 épocas
(100\% test). Versicolor vs. Virginica \textbf{no es separable}, con un techo de $\approx 83\%$
para un clasificador lineal. El MLP 4--8--3 alcanza consistentemente \textbf{91.1\% en test} para
todas las configuraciones de $\mu$ y $\eta$, con convergencia más rápida para $\mu=0.05$ (4
épocas al 95\% train). El momentum ($\eta=0.9$) no aporta beneficio en este problema pequeño y
bien condicionado, consistente con la observación del paper: ``el momentum por sí solo suele ser
de poco valor'' \cite{widrow1990}. La matriz de confusión concentra los 4 errores (de 45)
exclusivamente en el par versicolor $\leftrightarrow$ virginica.

El análisis de relevancia de variables (Tabla~\ref{tab:relevancia}) revela que los atributos de
pétalos concentran el 74.4\% del peso total de la primera capa (\textit{petal width}: 48.0\%,
\textit{petal length}: 26.4\%), consistente con el conocimiento botánico de que estas medidas
son los discriminantes más informativos en Iris \cite{fisher1936}.

\textbf{Two-moons (Figura~\ref{fig:moons}):} El MLP 2--8--1 con \bp~alcanza \textbf{93.3\% test},
trazando una frontera no lineal que se ajusta a la geometría de las medias lunas.

\textbf{Wine (Figura~\ref{fig:wine}):} Los clasificadores lineales alcanzan un rendimiento
sorprendentemente alto: Wiener = 98.1\%, \muLMS~= 98.1\%, \aLMS~= 96.3\%. Esto sugiere que
la clase 0 es mayoritariamente separable del resto en el espacio de 13 atributos fisicoquímicos.

\textbf{MNIST-04 (Figura~\ref{fig:mnist}):} El MLP 64--16--5 alcanza \textbf{98.5\% test},
demostrando que incluso una arquitectura simple con \bp~puede manejar datos de mediana
dimensionalidad (64 features) con alta precisión.

\subsection{E7: MRIII $\equiv$ Backpropagation}

\figura{fig_mriii_bp}{Izquierda: Error relativo $\norm{\nabla_{\text{MRIII}} -
\nabla_{\text{BP}}} / \norm{\nabla_{\text{BP}}}$ vs.~$\Delta s$ (escala log-log, eje
$x$ invertido). El error mínimo es $5.05 \times 10^{-8}$ en $\Delta s = 10^{-8}$.
Derecha: Costo computacional por presentación (ms).}{mriii_bp}

\tabla{tab_mriii}{Costo computacional de \bp~vs.~MRIII para redes 4--H--3.}{mriii}

La Figura~\ref{fig:mriii_bp}(a) confirma \textbf{H4} con precisión numérica: el gradiente
estimado por perturbación (MRIII) converge al gradiente analítico de \bp~con error relativo
mínimo de \textbf{$5.05 \times 10^{-8}$} en $\Delta s = 10^{-8}$. El error decrece linealmente
con $\Delta s$ hasta que la precisión de punto flotante de 64 bits ($\approx 10^{-16}$) causa
un rebote para $\Delta s < 10^{-8}$.

La Figura~\ref{fig:mriii_bp}(b) y la Tabla~\ref{tab:mriii} cuantifican el costo de esta
equivalencia: MRIII es $3.5\times$ más costoso para $H=4$ y \textbf{$74.1\times$} para $H=256$.
La razón crece aproximadamente como $\mathcal{O}(N_{\text{Adalines}})$, validando que \bp~es la
opción óptima para implementación digital, mientras MRIII tiene sentido únicamente en hardware
analógico donde la perturbación física del sumidero es más simple que la retropropagación de
señales de error.

\subsection{Tabla resumen}

\tabla{tab_resumen}{Resumen consolidado de todos los experimentos.}{resumen}

% ═══════════════════════════════════════════════════════════════════
%  5. CONCLUSIONES (20%)
% ═══════════════════════════════════════════════════════════════════
\section{Conclusiones}

\subsection{Aporte al conocimiento}

El artículo de Widrow \& Lehr (1990) realiza un aporte fundamental a la disciplina al
\textbf{unificar 30 años de algoritmos de aprendizaje supervisado} para RNA feedforward bajo un
principio común: la mínima perturbación. Los experimentos realizados en este trabajo verifican
empíricamente esta unificación:

\begin{enumerate}[leftmargin=*]
    \item \textbf{El linaje Perceptrón $\to$ Adaline $\to$ Madaline $\to$ Backpropagation} no es
    una colección arbitraria de algoritmos, sino una progresión lógica gobernada por el mismo
    principio geométrico: $\Delta W_k$ es siempre colineal con $X_k$, minimizando $\norm{\Delta
    W_k}$ para la corrección deseada.
    \item \textbf{El trade-off capacidad--generalización} ($C_s \approx 2N_w$, $N_p \gg N_w/N_y$)
    establece límites fundamentales que siguen siendo relevantes en la era del \textit{deep
    learning}: una red sobreparametrizada ($N_w > N_p$) memoriza; una red correctamente
    dimensionada generaliza.
    \item \textbf{El valor pedagógico e histórico} del artículo es inmenso: las ecuaciones de
    \bp~(1986) ya estaban implícitas en el trabajo de Werbos (1974), y la equivalencia
    MRIII $\equiv$ \bp~demuestra que el descubrimiento científico a menudo ocurre por múltiples
    caminos independientes.
\end{enumerate}

\subsection{Evaluación de las hipótesis}

\begin{description}
    \item[H1 -- Capacidad de Cover:] \textbf{Confirmada.} La simulación Monte Carlo (E1) reproduce
    la fórmula teórica con desviación media $\leq 0.034$.
    \item[H2 -- Convergencia del Perceptrón:] \textbf{Confirmada.} El Perceptrón converge en 2
    épocas en datos separables (E2), pero oscila $\pm5.5$ puntos en datos no separables (E3),
    mientras \muLMS~mantiene estabilidad ($\pm1.5$ pts, a 1.4 pts del óptimo de Wiener).
    \item[H3 -- Poder no lineal y óptimos locales:] \textbf{Confirmada con matices.} XOR es
    resoluble solo por arquitecturas multicapa, pero MRII (0/20) resultó más vulnerable a óptimos
    locales que MLP-\bp~(6/20 para 2--2--1, 11/20 para 2--3--1). El ruido esporádico no fue
    suficiente para rescatar MRII en nuestra implementación.
    \item[H4 -- Equivalencia MRIII $\equiv$ Backpropagation:] \textbf{Confirmada.} El error
    relativo entre gradientes es $5.05 \times 10^{-8}$, y el costo de MRIII escala como
    $\mathcal{O}(N_{\text{Adalines}})$, siendo $3.5\times$ a $74.1\times$ más costoso que \bp.
\end{description}

\subsection{Limitaciones}

\begin{itemize}[leftmargin=*]
    \item Los datasets utilizados son de baja dimensionalidad ($\leq 64$ atributos). No se evaluó
    la escalabilidad a problemas de alta dimensionalidad.
    \item La implementación de MRII, aunque conceptualmente completa, no logró resolver XOR ni
    siquiera con el mecanismo de ruido. Una implementación más refinada del escape (recocido
    simulado, \textit{random restarts}) podría mejorar estos resultados.
    \item No se implementaron variantes avanzadas de \bp~mencionadas en el artículo
    (\textit{quickprop}, \textit{delta-bar-delta}, \textit{backprop through time}).
    \item Las comparaciones se limitan a los algoritmos del artículo; no se incluyen métodos
    modernos (SVM, Random Forest, \textit{transformers}) que podrían establecer una línea base
    más exigente.
\end{itemize}

\subsection{Reflexión personal}

Treinta y cinco años después de su publicación, el artículo de Widrow \& Lehr permanece
sorprendentemente vigente. El principio de mínima perturbación es, en esencia, el mismo que
gobierna el \textit{stochastic gradient descent} (SGD) que entrena los modelos de \textit{deep
learning} actuales: $\mu$-\textit{LMS} generalizado a arquitecturas profundas. El momentum
(Eqs.~96--97) sigue siendo un hiperparámetro estándar en optimizadores como Adam. Y la
perturbación de MRIII anticipa los métodos de optimización de orden cero (\textit{zeroth-order
optimization}) usados hoy en escenarios donde el gradiente no está disponible. La gran lección
del artículo es que los fundamentos importan: comprender la geometría del aprendizaje —el
hiperplano separador, el paraboloide de MSE, las mesetas del error signum— es esencial para
diseñar, depurar y mejorar cualquier sistema de aprendizaje automático, clásico o moderno.

% ═══════════════════════════════════════════════════════════════════
%  REFERENCIAS (10%)
% ═══════════════════════════════════════════════════════════════════
\begin{thebibliography}{99}

\bibitem{widrow1990}
B.~Widrow and M.~A.~Lehr,
``30 years of adaptive neural networks: Perceptron, Madaline, and backpropagation,''
\textit{Proceedings of the IEEE}, vol.~78, no.~9, pp.~1415--1442, 1990.

\bibitem{rosenblatt1958}
F.~Rosenblatt,
``The perceptron: A probabilistic model for information storage and organization in the brain,''
\textit{Psychological Review}, vol.~65, no.~6, pp.~386--408, 1958.

\bibitem{widrow1960}
B.~Widrow and M.~E.~Hoff,
``Adaptive switching circuits,''
in \textit{IRE WESCON Convention Record}, vol.~4, 1960, pp.~96--104.

\bibitem{cover1965}
T.~M.~Cover,
``Geometrical and statistical properties of systems of linear inequalities with applications
in pattern recognition,''
\textit{IEEE Transactions on Electronic Computers}, vol.~EC-14, no.~3, pp.~326--334, 1965.

\bibitem{rumelhart1986}
D.~E.~Rumelhart, G.~E.~Hinton, and R.~J.~Williams,
``Learning representations by back-propagating errors,''
\textit{Nature}, vol.~323, pp.~533--536, 1986.

\bibitem{werbos1974}
P.~J.~Werbos,
``Beyond regression: New tools for prediction and analysis in the behavioral sciences,''
Ph.D.~dissertation, Harvard University, Cambridge, MA, 1974.

\bibitem{minsky1969}
M.~L.~Minsky and S.~A.~Papert,
\textit{Perceptrons: An Introduction to Computational Geometry}.
Cambridge, MA: MIT Press, 1969.

\bibitem{cybenko1989}
G.~Cybenko,
``Approximation by superpositions of a sigmoidal function,''
\textit{Mathematics of Control, Signals and Systems}, vol.~2, no.~4, pp.~303--314, 1989.

\bibitem{fisher1936}
R.~A.~Fisher,
``The use of multiple measurements in taxonomic problems,''
\textit{Annals of Eugenics}, vol.~7, no.~2, pp.~179--188, 1936.

\bibitem{widrow1987}
B.~Widrow, R.~G.~Winter, and R.~A.~Baxter,
``Layered neural nets for pattern recognition,''
\textit{IEEE Transactions on Acoustics, Speech, and Signal Processing}, vol.~36, no.~7,
pp.~1109--1118, 1988.

\bibitem{sejnowski1987}
T.~J.~Sejnowski and C.~R.~Rosenberg,
``Parallel networks that learn to pronounce English text,''
\textit{Complex Systems}, vol.~1, pp.~145--168, 1987.

\bibitem{nguyen1989}
D.~H.~Nguyen and B.~Widrow,
``Neural networks for self-learning control systems,''
\textit{IEEE Control Systems Magazine}, vol.~10, no.~3, pp.~18--23, 1990.

\bibitem{hastie2009}
T.~Hastie, R.~Tibshirani, and J.~Friedman,
\textit{The Elements of Statistical Learning: Data Mining, Inference, and Prediction},
2nd ed. New York: Springer, 2009.

\bibitem{bishop2006}
C.~M.~Bishop,
\textit{Pattern Recognition and Machine Learning}.
New York: Springer, 2006.

\end{thebibliography}

\end{document}
